\chapter{Estado del arte del problema a abordar}

El uso de Octrees como estructura de datos en la gestión y almacenamiento de nubes de puntos 3D es en la actualidad predominante a nivel global. Existen librerías de software libre ampliamente reconocidas que proporcionan soporte nativo tanto a estructuras Octrees como a KD-Trees para optimizar la indexación espacial.

Una de estas librerías es Open3D. Está escrita en C++ pero cuenta con interfaces estables para numerosos lenguajes, considerándose de facto el estándar en Python para la manipulación y visualización 3D. Open3D incluye métodos avanzados para el cálculo de distancias y la búsqueda de vecindad, dando soporte tanto a consultas basadas en kernels de radio como a la técnica de los $K$ vecinos más cercanos (\textit{$K$-Nearest Neighbors} o KNN).

Otro gigante en el manejo de nubes de puntos, considerada la librería de referencia absoluta en el entorno académico e industrial para el procesamiento de geometría 3D en C++, es PCL (\textit{Point Cloud Library}) [ref pcl]. PCL proporciona un ecosistema masivo y modular que da soporte nativo a estructuras espaciales mediante las clases \texttt{pcl::Octree} y \texttt{pcl::KdTree}. El enfoque de PCL en sus Octrees tradicionales se basa en una estructura dinámica de punteros. Este diseño otorga una gran versatilidad para realizar modificaciones en tiempo real sobre la nube, como la inserción o eliminación dinámica de puntos y la detección de cambios espaciales. Por la contra, introduce una severa penalización en el rendimiento de la memoria. Esto se debe a la incoherencia entre la localidad espacial de los datos y su disposición física en el \textit{heap} o montón, provocando constantes fallos de caché y saltos de punteros en el direccionamiento de la CPU durante los recorridos de búsqueda.

Un aspecto crítico que comparten tanto Open3D como PCL es que no implementan sus algoritmos de búsqueda desde cero, sino que integran internamente la librería FLANN (\textit{Fast Library for Approximate Nearest Neighbors}) [ref flann]. FLANN es una biblioteca altamente optimizada para la búsqueda de vecinos más cercanos en espacios vectoriales multidimensionales. Si bien su principal característica es la capacidad de realizar búsquedas aproximadas (\textit{Approximate Nearest Neighbor} o ANN) para mitigar la maldición de la dimensionalidad sacrificando un margen mínimo de precisión, en el contexto de búsquedas exactas por radio o kernel (como las evaluadas en este estudio) se configura en su modo exacto para operar como línea base de comparación directa.

Otras librerías relevantes con el mismo propósito, escritas también en C++, son PicoTree y NanoTree. PicoTree genera un KD-Tree integrado en un array contiguo mediante la indexación previa de los puntos, permitiendo búsquedas esféricas y de KNN altamente eficientes. Por su parte, NanoTree destaca por ser una implementación ligera y compacta optimizada para sistemas con restricciones de recursos.

Por otro lado, las curvas de llenado del espacio (SFC) son ampliamente aplicadas a estructuras jerárquicas para tareas de indexado o compresión. La integración de códigos de Morton en estructuras de Octree lineales ha sido una estrategia clave para mejorar la eficiencia en la ordenación de puntos. Un ejemplo de ello es el trabajo de Behley et al. con la librería \textit{unibn} [2], donde se emplean estos códigos para reducir la necesidad de recorrer el árbol en toda su profundidad durante las búsquedas de vecindad, logrando superar el rendimiento de otras herramientas populares de la disciplina. Estos códigos son reconocidos por su simplicidad algorítmica y rapidez de cómputo a nivel de bits. Sin embargo, el otro gran algoritmo de SFC, las curvas de Hilbert, ha ganado una notable popularidad en la literatura reciente por garantizar una preservación superior de la localidad espacial, asegurando que si dos puntos son cercanos en el espacio tridimensional, mantendrán dicha proximidad en la proyección unidimensional de memoria de forma estricta.

No obstante, la aplicación directa de SFC en la búsqueda de vecindades en nubes de puntos 3D obtenidas con sensores LiDAR fue implementada originalmente en el trabajo de Díaz Viñambres et al. [1]. En dicha investigación se analiza en primer lugar una estructura Octree de punteros estándar construida mediante fragmentaciones recursivas. A pesar de que a esta estructura se le pueden aplicar reordenamientos según curvas de Morton o Hilbert para mejorar la cercanía en memoria de los puntos próximos, el enfoque clásico sigue adoleciendo del problema de los accesos indexados: cada transición entre nodos durante el recorrido implica saltar a una dirección del mapa de memoria que puede estar totalmente dispersa.

Por este motivo, los autores propusieron e implementaron una versión lineal del Octree, la cual almacena todos los puntos de la nube en un único vector contiguo, garantizando que todos los elementos pertenecientes a un mismo nodo hoja residan en posiciones físicas consecutivas del array. Adicionalmente, se introdujo un algoritmo de búsqueda por kernel geométrico que optimiza el proceso de recorrido: a diferencia del método clásico que desciende recursivamente verificando la mera intersección, este algoritmo identifica de forma temprana si el volumen del kernel contiene por completo el área que cubre el nodo (sea este interno o terminal). De cumplirse esta condición de inclusión total, el algoritmo realiza un volcado secuencial directo de todos los puntos del subárbol, abortando el descenso y acelerando drásticamente la consulta.

En la publicación de Díaz Viñambres et al. [1] se presentan resultados experimentales que demuestran la superioridad del Octree lineal (en sus variantes con reordenamiento de Morton y Hilbert) frente al Octree clásico de punteros. Asimismo, se incluye un exhaustivo análisis comparativo de rendimiento frente a las soluciones de código abierto de referencia mencionadas anteriormente, tales como los módulos de PCL, PicoTree, NanoTree y la librería \textit{unibn}. Bajo diferentes tipologías de datasets y variaciones en el radio del kernel de búsqueda, el Octree lineal reportó, de manera sistemática, los menores tiempos de ejecución computacional.

Sobre este escenario técnico se fundamenta el presente trabajo. En la práctica totalidad de las librerías del estado del arte existe un patrón de diseño común: las estructuras justifican su eficiencia basándose exclusivamente en la aplicación de la poda a nivel de nodos intermedios (ramas). Mediante funciones de solapamiento geométrico, se descartan de forma temprana subárboles completos cuyos volúmenes no intersecan con el kernel de consulta. Sin embargo, una vez que el algoritmo alcanza un nodo hoja válido, se ve obligado a evaluar la distancia de todos los puntos contenidos en él de manera individual y secuencial.

Hasta donde se tiene constancia en la literatura técnica y en las soluciones de código abierto analizadas, no existe ninguna implementación que acote el rango de puntos a evaluar mediante una poda matemática interna que actúe de forma nativa \textbf{en el interior del propio nodo hoja}.

Este cuello de botella constituye la principal motivación y el núcleo de la contribución de este trabajo. Al introducir funciones de cálculo de rango basadas en reordenamientos geométricos locales (tanto en sistemas cartesianos como esféricos y cilíndricos), este estudio permite aislar zonas y subrangos de puntos dentro de los nodos hoja. Esto cambia la concepción que se tenía sobre los nodos hoja, anteriormente pensados como unidades atómicas e indivisibles. Este enfoque persigue un incremento en el rendimiento temporal de la estructura del Octree lineal, motivada por la reducción del número de puntos evaluados como candidatos a vecinos. Como un extra, también busca aportar una mayor flexibilidad al diseño del software, permitiendo configurar tamaños de hoja más elevados sin penalizar el tiempo de cómputo por desbordamiento de búsquedas exhaustivas.




